\section{Introduction}
\label{sec:intro}

Survey research is a foundational instrument across the social, behavioral, and marketing sciences, yet it faces rising costs, declining response rates, and a growing threat from automated responses that contaminate online panels~\citep{westwood2025potential,camatarri2024predicting}. Against this backdrop, \emph{silicon sampling}---the use of large language models (LLMs) to stand in for human respondents---has emerged as an attractive method for augmenting or partially substituting costly data collection~\citep{sarstedt2024using,anthis2025social,bail2024generative}. The most consequential and most difficult variant of this task is \emph{individual-level} silicon sampling: rather than reproducing an aggregate marginal for a subpopulation, the goal is to predict how a \emph{specific} respondent, characterized by a demographic profile and a history of prior answers, will respond to a \emph{specific} new question. Individual-level prediction is what enables counterfactual reasoning, panel imputation, and personalized instrument design, and it is precisely where current methods are weakest.

The dominant paradigm conditions an LLM on a persona and reads its generated answer token as the prediction~\citep{sun2024random,hu2024quantifying,li2026personabased}. This paradigm suffers from three well-documented pathologies. First, it is unreliable and systematically biased: model answers diverge from human reference distributions in ways that persona conditioning does not fix~\citep{dominguezolmedo2023questioning,morocho2026assessing,taubenfeld2024systematic}. Second, it is poorly calibrated, because a single sampled token conveys no faithful uncertainty over the option set~\citep{geng2024survey}. Third, and most damaging for our task, it transfers poorly: a prompt tuned for one instrument, domain, or time period rarely holds up on unseen items, other domains, or later survey waves~\citep{ku2026silicon}. We contend that these are not prompt-engineering nuisances but symptoms of a structural mistake---forcing the model to \emph{commit} to a discrete answer collapses the rich latent structure of how a respondent would react into a single lossy token.

Our starting point is an analogy that we develop rigorously. In spatio-temporal learning, \citet{he2025geolocation} showed that an LLM is far more useful as a \emph{representation generator} than as a direct predictor: instead of asking the LLM to forecast demand at a location, they extract an LLM-derived geolocation vector and feed it to a lightweight downstream model, obtaining a generic enhancer that transfers across cities. We transpose this ``predictor $\rightarrow$ representation generator'' move from the spatial domain to individual behavior prediction. Rather than sampling an answer token, we use the LLM to produce a latent \emph{reaction} vector for a respondent facing a question, and delegate the actual choice to a small, calibratable head. Crucially, we extract this vector with a \emph{generative} LLM embedder~\citep{behnamghader2026llmvecgen} that preserves the model's output semantics, rather than a purely discriminative encoder~\citep{behnamghader2024llmvec} that discards them, because the object we wish to represent---a latent disposition to answer---is generative in nature. This distinguishes our representation from raw hidden states, which are geometrically anisotropic and uncalibrated, and from generic text embeddings, which are not aligned to option-level choice.

Concretely, we introduce \textbf{ResponseVec}, a query-conditioned latent reaction vector, and its reusable counterpart \textbf{RespondentVec}. A literature-grounded survey prompt fuses the respondent's demographic profile, their target-relevant answer history retrieved with option information intact, and the current item with its full option set; the generative embedder maps this prompt to ResponseVec; and an \emph{option-aware decoder} scores each option by its interaction with the reaction vector and is calibrated on a small held-out set. RespondentVec instead encodes only the profile and history once, injecting the item at the decoder, so that a single vector is reused across an entire questionnaire. This pairing exposes a fundamental tension---query conditioning maximizes task adaptivity at high per-item cost, whereas reuse amortizes compute at the price of item specificity---and we characterize the resulting trade-off explicitly.

We evaluate under three transfer regimes that jointly stress the paradigm: \emph{unseen questions} (held-out items, same respondents), \emph{cross-domain} transfer (train on three ISSP domains, test on the held-out domain), and \emph{out-of-distribution demographic intersections} (hold out entire demographic-combination cells of respondents), the last of which directly targets the transfer-evaluation gap identified by \citet{ku2026silicon}. Against seven baseline families---direct generation with progressive calibration, classical psychometric models (majority and MIRT), item-conditional collaborative filtering, quantized fine-tuning~\citep{mao2024survey,song2025enhancing}, generic text vectors~\citep{wang2024multilingual}, raw hidden states, and a non-generative LLM2Vec ablation---we test accuracy, calibration, robustness, transfer, and compute.

This paper makes four contributions. \textbf{(i) Conceptual:} we are the first to transpose the ``predictor $\rightarrow$ representation generator'' paradigm from spatio-temporal learning to individual behavior prediction, and we formalize a dual ResponseVec/RespondentVec representation framework. \textbf{(ii) Methodological:} we design an option-aware, calibratable decoder that anchors a lightweight head to the option distribution of a concrete instrument, and we justify the use of a generative embedder to carry latent reaction semantics. \textbf{(iii) Empirical:} we provide the first systematic comparison of seven baseline families under a unified three-way transfer protocol, filling the OOD-demographic evaluation gap. \textbf{(iv) Analytical:} we chart the accuracy--transfer--compute boundary between query-conditioned and reusable representations, yielding operational guidance for practitioners. The remainder of the paper reviews related work, formalizes the task, details the method and the trade-off, and reports our experimental design and analysis.
